%! The Statistical Cycle and Types of Data — Unit 1, Lesson 1.1 — Guided Notes & Practice (Answer Key)
% THE in-class document, in the MAIN IDEAS / NOTES shape (LESSON_SHAPE.md §1, ported from AP
% Statistics 2026-09-12): the vocabulary box, the hook, then ONE two-column guidednotes table.
% Each row is one idea — a label on the left; on the right one or two COMPLETE printed
% sentences, ONE large pre-drawn display the student reads or names, then a bold prompt and a
% two-across grid of numbered problems with 1.2–2.4cm of writing room. Problems run 1..19.
% DENSITY: a \blank{} only where a word or number IS the answer (the frequency table in row 3,
% problem 9 — the one \wordbank strip sits above it); no mid-sentence blanks; every display is
% pre-drawn; the page belongs to the student's pen.
%   I DO   : the teacher reads the definition, marks up the display, works the first problem
%            of each grid (1, 5, 9's bus row, 12).
%   WE DO  : the class works the rest of each grid, students holding the pen; the last row,
%            Guided Practice (16–19), is worked entirely together on the council's SECOND
%            survey (the cycle turned).
%   YOU DO : nothing on this page — ../homework/ IS the individual practice, started in the
%            last ten minutes of the period. No independent set, no reflection box, no
%            objectives box (the cover carries the targets).
% THE TRAP is problem 8 (a question that passes all three tests as written — the student who
% "repairs" it has learned to rewrite, not to test).
% THE CRUX is problem 14 (row 4, "The Question Decides"): the council's arithmetic is correct
% (6 of 60 = 10%, about 90 of 900) and the sentence "about 90 students WOULD bike if there
% were racks" is still unsupported — the question asked how students travel NOW, so the
% problem entered at stage 1. The hook vote is re-taken there (PS.DC.1a–c). Problem 15 is the
% breakfast-cart counterweight (a stage-2 entry); problem 19 is the crux transferred.
%
% THE DATA
%   Riverbend Student Council, last spring: 60 students asked how they usually get to school —
%     bus 27 (45%), car 15 (25%), walk 12 (20%), bike 6 (10%); 900 students; 0.10 x 900 = 90
%   Breakfast cart (problem 15): 40 students leaving the cart, 14 bought a bagel -> 35%
%   Guided Practice, the council's second survey: 80 asked "would you ride a bike to school if
%     there were racks?", 24 yes -> 30%; 0.30 x 900 = 270
%
% PAGE LOCKSTEP with notes_key/: the key is GENERATED from this file by templates/lesson/mkkey.py
% (spec: templates/lesson/examples/lesson01_notes_key.json) — every \blank{} becomes an \ans{},
% every \vterm{} a \vtermans{}{}, every \par\writelines{n} n \ansline{} calls, and the answer
% argument of every \pcell, \writespace and \labelbox is filled. Answer spaces are fixed
% heights, so the two files cannot drift. KEEP EVERY \pcell ON ONE SOURCE LINE — mkkey matches
% line by line.
% TABLE MECHANICS: inside a cell, `\\` ENDS THE ROW — break lines with \par. A row cannot break
% across pages: the figure gets its own sub-row (`\\` then `& ...`) and EACH GRID ROW is its own
% sub-row — close the probgrid, `\\ &`, reopen as probgrid* (no top rule).
\documentclass[12pt]{article}
\usepackage{saar-article}
\usepackage{saar-key}   % pulls in -boxes; do NOT also load -boxes
\newcommand{\TallMath}[1]{$\displaystyle #1\rule[-1.4em]{0pt}{3.2em}$}
% Word bank strip — ONLY above a table the student fills (row 3, problem 9). Defined per
% document; the key inherits it and prints the same strip, so heights match.
\newcommand{\wordbank}[1]{\par\vspace{2pt}\noindent\fcolorbox{goldacc}{goldbg}{\parbox{\dimexpr\linewidth-2\fboxsep-2\fboxrule\relax}{\small\textbf{Word bank:} #1}}\par\vspace{4pt}}
% Vocabulary rows: a FIXED-HEIGHT row carrying the term label and OPEN writing space —
% no underline beside the term and no rule line (user correction 2026-09-06: the black
% answer line crowded the row; students write beside and below the term). The key fills
% exactly the same height, so the box cannot drift blank vs. keyed. saar-article's
% \termblank adds an inline blank and a rule line, so the pair is defined here (the key is
% generated from this file and inherits both). 2.3cm per row gives students three
% handwritten lines of room; keep key definitions to two lines.
\newcommand{\termrowheight}{2.3cm}
\newlength{\termrowinset}\setlength{\termrowinset}{7pt}
\newcommand{\vterm}[1]{%
  \par\noindent\begin{minipage}[t][\termrowheight][t]{\linewidth}%
    \vspace*{\termrowinset}%
    \textbf{\textcolor{cerulean}{#1:}}%
  \end{minipage}\par}
\newcommand{\vtermans}[2]{%
  \par\noindent\begin{minipage}[t][\termrowheight][t]{\linewidth}%
    \vspace*{\termrowinset}%
    \textbf{\textcolor{cerulean}{#1:}}\quad\ans{#2}%
  \end{minipage}\par}

\begin{document}

\pageheader{Unit 1, Lesson 1.1}{Guided Notes \& Practice --- Answer Key}
% NAMESTRIP: no \namedateperiod here — mirrors the blank. NO teachernote here —
% teacher prose lives in the lesson plan.

\vspace{0.15in}

\begin{vocabbox}
\small
Fill in each term \textbf{as we name it} in the notes below.
\vtermans{Data cycle}{The four repeating stages of a study --- formulate, collect, organize, analyze and communicate. Each stage can use only what the stage before it handed over.}
\vtermans{Statistical question}{A question whose answers vary from one individual to the next, that names the group it is about, and that says what will be recorded for each one.}
\vtermans{Variable of interest}{The characteristic the question says you will record for each individual. The question names it before any data exist.}
\vtermans{Relative frequency}{A count divided by the total, written as a decimal or a percent --- the share of the whole that one category makes up.}
\end{vocabbox}

\boxguard[12]
\begin{hookbox}
\small
The Riverbend Student Council wants \textbf{more bike racks}. Last spring they asked $60$
students how they usually get to school; $6$ said \emph{bike}. Their claim: \textit{``$10\%$ bike
now, so a lot more would if there were somewhere to park.''} The arithmetic is right.

\vspace{3pt}
\textbf{Does the survey support the claim?} Circle one: \quad yes \qquad no \qquad can't tell
\quad --- and say why you think so. We vote again at problem~14.
\par\ansline{Answers vary --- many vote yes because the 10\% checks out. (Settled at problem 14: the}
\ansline{question asked how students travel NOW, so the claim entered at stage 1.)}
\end{hookbox}

\small
\begin{guidednotes}
% ── ROW 1: the four stages ───────────────────────────────────────────────────
\mainidea[Four stages, in order]{Data Cycle} &
Every study runs the same four stages, in order, and then starts over: the \textbf{data
cycle}. \textbf{Each stage can only use what the stage before it handed over} --- if stage 1
never names the group, stage 4 cannot say who the conclusion is about.
\\
& {\centering
\begin{tikzpicture}
  % four stage boxes in a loop (real cm; the Notes column is about 12cm wide)
  \foreach \i/\num/\sname/\hands in {
      0/1/{Formulate}/{a question naming\\a group and a\\measurement},
      1/2/{Collect}/{the raw values,\\one row per\\individual},
      2/3/{Organize}/{a table that\\makes the\\pattern visible},
      3/4/{Analyze \&\\communicate}/{a conclusion in\\context --- and\\the next question}} {
    \fill[frost, rounded corners=5pt] ({\i*2.85},0) rectangle ({\i*2.85+2.6},2.05);
    \draw[cerulean, thick, rounded corners=5pt] ({\i*2.85},0) rectangle ({\i*2.85+2.6},2.05);
    \node[circle, fill=cerulean, text=white, font=\bfseries\scriptsize, inner sep=1.5pt]
      at ({\i*2.85+0.3},1.78) {\num};
    \node[font=\footnotesize\bfseries\color{cerulean}, align=center, anchor=north]
      at ({\i*2.85+1.4},1.98) {\sname};
    \node[font=\tiny, align=center] at ({\i*2.85+1.3},0.55) {\hands};
  }
  \foreach \i in {0,1,2} { \draw[->, thick, charcoal] ({\i*2.85+2.62},1.0) -- ({\i*2.85+2.83},1.0); }
  % the loop back to stage 1
  \draw[->, very thick, goldacc] (9.85,-0.05) -- (9.85,-0.42) -- (1.3,-0.42) -- (1.3,-0.05);
  \node[font=\tiny\itshape\color{goldacc}, fill=white, inner sep=1pt] at (5.6,-0.42) {the answer raises the next question --- and the cycle turns};
\end{tikzpicture}\par\vspace{4pt}
The stage that names the group a conclusion may describe: \labelbox{2.6cm}{stage 1}\par}
\\
& \notesprompt{Here is what the council did last spring, out of order. Give the stage number, and say what that stage handed on.}
\begin{probgrid}{|Y|Y|}
\pcell{1}{Asked $60$ students how they usually get to school.}{1.2cm}{Stage 2 --- collect. Hands on 60 raw answers.} & \pcell{2}{Told the district what the survey showed.}{1.2cm}{Stage 4 --- analyze. Hands on a conclusion.} \\ \hline
\end{probgrid}
\\
& \begin{probgrid}*{|Y|Y|}
\pcell{3}{Decided to ask how Riverbend students travel to school.}{1.2cm}{Stage 1 --- formulate. Hands on the question.} & \pcell{4}{Counted the $60$ answers into a table by travel method.}{1.2cm}{Stage 3 --- organize. Hands on the table.} \\ \hline
\end{probgrid}
\\ \hline
% ── ROW 2: stage 1 — the three tests ─────────────────────────────────────────
\mainidea[Stage 1]{Three Tests} &
A \textbf{statistical question} passes three tests: the answers \textbf{vary} from one
individual to the next, it names the \textbf{group} it is about, and it says what will be
\textbf{recorded} for each one. \textbf{The thing you record is the variable of interest.}
\\
& {\centering
\begin{tikzpicture}
  % the council's question passing through three gates (real cm)
  \node[font=\footnotesize\itshape, anchor=west] at (0,2.75)
    {``How do \textbf{Riverbend students} usually \textbf{get to school}?''};
  \draw[->, thick, charcoal] (0.6,2.5) -- (0.6,2.2);
  \foreach \i/\test/\ask in {
      0/{1 \; VARY?}/{would two students\\answer differently?},
      1/{2 \; GROUP?}/{who is it about?},
      2/{3 \; RECORD?}/{for each one I will\\write down \ldots}} {
    \fill[goldbg, rounded corners=4pt] ({\i*3.3},0.2) rectangle ({\i*3.3+3.0},2.0);
    \draw[goldacc, thick, rounded corners=4pt] ({\i*3.3},0.2) rectangle ({\i*3.3+3.0},2.0);
    \node[font=\small\bfseries\color{goldacc}] at ({\i*3.3+1.5},1.6) {\test};
    \node[font=\tiny, align=center] at ({\i*3.3+1.5},0.85) {\ask};
  }
  \draw[->, thick, charcoal] (3.02,1.1) -- (3.28,1.1);
  \draw[->, thick, charcoal] (6.32,1.1) -- (6.58,1.1);
  \draw[->, very thick, greenacc] (9.65,1.1) -- (10.2,1.1);
  \node[font=\footnotesize\bfseries\color{greenacc}, anchor=west] at (10.2,1.1) {ready};
\end{tikzpicture}\par\vspace{4pt}
For this question the variable of interest is \labelbox{4.2cm}{how a student travels}\par}
\\
& \notesprompt{Which test does each question fail? Repair it --- keep the original subject and put back only what was missing.}
\begin{probgrid}{|Y|Y|}
\pcell{5}{``How do I get to school?''}{1.4cm}{Fails \emph{vary}. Repair: ``How do Riverbend students get to school?''} & \pcell{6}{``Do people like biking?''}{1.4cm}{Fails \emph{group}. Repair: ``What percent of Riverbend students bike?''} \\ \hline
\end{probgrid}
\\
& \begin{probgrid}*{|Y|Y|}
\pcell{7}{``Is the breakfast cart good?''}{1.4cm}{Fails \emph{record} --- ``good'' is not something you write down.} & \pcell{8}{``How many minutes is a Riverbend student's trip to school?'' \textbf{Test it before you touch it.}}{1.4cm}{It passes all three. Nothing to repair.} \\ \hline
\end{probgrid}
\\ \hline
% ── ROW 3: stages 2 and 3 — the frequency table ──────────────────────────────
\mainidea[Stages 2 and 3]{Frequency Table} &
Stage 2 collects one answer per individual. Stage 3 organizes them into a \textbf{frequency
table}: each category, its \textbf{frequency} (the count), and its \textbf{relative frequency}
(count $\div$ total, as a percent). \textbf{Stage 3 adds nothing new --- it makes the pattern visible.}
\\
& \textbf{9.} The council's $60$ answers --- fill in the missing percents and the totals.
\wordbank{$15/60$ \;\textbullet\; $12/60$ \;\textbullet\; $6/60$ \;\textbullet\; $0.25$ \;\textbullet\; $0.20$ \;\textbullet\; $0.10$ \;\textbullet\; $25\%$ \;\textbullet\; $20\%$ \;\textbullet\; $10\%$ \;\textbullet\; $60$ \;\textbullet\; $100\%$}
\footnotesize\renewcommand{\arraystretch}{1.2}
\begin{tabularx}{\linewidth}{@{} l c X @{}}
\toprule
\textbf{How they get to school} & \textbf{Frequency (count)} & \textbf{Relative frequency (percent)} \\
\midrule
Bus  & $27$ & $27/60 = 0.45 = 45\%$ \\
Car  & $15$ & \ans{$15/60 = 0.25 = 25\%$} \\
Walk & $12$ & \ans{$12/60 = 0.20 = 20\%$} \\
Bike & $6$  & \ans{$6/60 = 0.10 = 10\%$} \\
\midrule
\textbf{Total} & \ans{$60$} & \ans{$60/60 = 1.00 = 100\%$} \\
\bottomrule
\end{tabularx}\small
\\
& \notesprompt{Read the table --- then move to stage 4.}
\begin{probgrid}{|Y|Y|}
\pcell{10}{Add the four percents. What must they total? Which group is smallest --- and could you see that in the $60$ raw answers?}{1.6cm}{$100\%$. Bike is smallest ($10\%$) --- invisible in 60 raw answers.} & \pcell{11}{Riverbend has $900$ students. Scale the bike percent to the school, then write the council's \textbf{conclusion in context}: the group asked, the number, what was measured.}{1.8cm}{$0.10 \times 900 = 90$. ``Of the 60 asked, 10\% bike --- about 90 of our 900.''} \\ \hline
\end{probgrid}
\\ \hline
% ── ROW 4: THE CRUX — the question decides the type AND the claim ────────────
\mainidea[Stage 1 rules the rest]{The Question Decides} &
\textbf{The question decides the type} --- before any data exist --- and the type decides
whether you summarize with a mean or with percents. The question also decides what stage 4
may claim: \textbf{a conclusion answers only the question that was asked, about only the
group that was asked.}
\\
& {\centering
\begin{tikzpicture}
  % the same 30 students, two questions, two types (real cm)
  \fill[frost, rounded corners=6pt] (0,0) rectangle (2.7,2.6);
  \draw[cerulean, thick, rounded corners=6pt] (0,0) rectangle (2.7,2.6);
  \node[font=\footnotesize\bfseries\color{cerulean}, align=center] at (1.35,2.15) {the same\\$30$ students};
  \foreach \i in {0,...,5} { \foreach \j in {0,...,4} {
    \fill[slate] ({0.45+\i*0.36},{0.35+\j*0.28}) circle (0.06); }}
  % question A
  \draw[->, thick, charcoal] (2.8,1.95) -- (3.5,1.95);
  \fill[goldbg, rounded corners=4pt] (3.6,1.4) rectangle (11.7,2.55);
  \draw[goldacc, thick, rounded corners=4pt] (3.6,1.4) rectangle (11.7,2.55);
  \node[font=\scriptsize\itshape, anchor=west] at (3.7,2.25) {``How many mornings a week do you eat breakfast at school?''};
  \node[font=\scriptsize, anchor=west] at (3.7,1.72) {record a \textbf{count} of mornings \; $\rightarrow$ \; \textbf{quantitative} \; $\rightarrow$ \; a mean};
  % question B
  \draw[->, thick, charcoal] (2.8,0.65) -- (3.5,0.65);
  \fill[frost, rounded corners=4pt] (3.6,0.05) rectangle (11.7,1.2);
  \draw[cerulean, thick, rounded corners=4pt] (3.6,0.05) rectangle (11.7,1.2);
  \node[font=\scriptsize\itshape, anchor=west] at (3.7,0.92) {``Do you ever eat breakfast at school?''};
  \node[font=\scriptsize, anchor=west] at (3.7,0.38) {record \textbf{yes or no} \; $\rightarrow$ \; \textbf{categorical} \; $\rightarrow$ \; percents};
\end{tikzpicture}\par\vspace{4pt}
Same students, two types. What decided the type? \labelbox{2.6cm}{the question}\par}
\\
& \textbf{In your own words:} the council's arithmetic is correct. Why can a correct percent
still fail to support the sentence the council wants to say?
\writespace{1.2cm}{The percent is right, but it answers a different question --- how students travel now, not what they would do.}
\\
& \notesprompt{Type first, from the question alone --- then what the question allows.}
\begin{probgrid}{|Y|Y|}
\pcell{12}{``How old are the bikes Riverbend students ride to school?'' Nothing is collected yet. What gets recorded, what type, summarized how?}{1.4cm}{The bike's age in years. Quantitative --- use a mean.} & \pcell{13}{``How do Riverbend students usually get to school?'' Why is a \emph{mean travel method} impossible --- and what does the table give instead?}{1.4cm}{The answers are labels; no arithmetic on them. Counts and percents instead.} \\ \hline
\end{probgrid}
\\
& \begin{probgrid}*{|Y|Y|}
\pcell{14}{\textbf{Vote again:} yes \; no \; can't tell. The council's arithmetic is right: $6$ of $60$ is $10\%$, about $90$ of $900$. Does the survey show students \emph{would} bike if there were racks? Which stage did the problem enter at?}{1.8cm}{No --- the question asked how they travel NOW, not what they would do. Entered at stage 1.} & \pcell{15}{Of $40$ students leaving the breakfast cart, $14$ bought a bagel --- $35\%$, correctly. The newspaper prints: \emph{``$35\%$ of Riverbend students buy a bagel every morning.''} Which stage did that enter at? Name \emph{two} claims the data cannot support.}{1.8cm}{Stage 2. (1) Only cart customers, not Riverbend students. (2) One morning, not every morning.} \\ \hline
\end{probgrid}
\\ \hline
% ── ROW 5: GUIDED PRACTICE — we do, together: the cycle turns ────────────────
\mainidea[Guided practice]{The Cycle Turns} &
The council goes back to stage 1 with the question the first survey could not answer:
\emph{``Would you ride a bike to school if there were racks?''} They ask $80$ Riverbend
students; $24$ say \emph{yes}.
\\
& {\centering
\begin{tikzpicture}[scale=0.88]
  % the second survey, drawn in real cm
  \fill[frost, rounded corners=6pt] (0,0) rectangle (11.4,2.25);
  \draw[cerulean, thick, rounded corners=6pt] (0,0) rectangle (11.4,2.25);
  \node[font=\small\bfseries\color{cerulean}] at (5.7,1.95) {Riverbend High School --- $900$ students};
  \fill[goldbg] (2.7,0.85) ellipse (2.45 and 0.66);
  \draw[goldacc, very thick] (2.7,0.85) ellipse (2.45 and 0.66);
  \node[font=\scriptsize\bfseries] at (2.7,1.05) {$80$ asked ``would you ride?''};
  \node[font=\scriptsize] at (2.7,0.65) {$24$ yes \quad $56$ no};
  \node[font=\scriptsize, anchor=west] at (5.5,1.2) {\textbf{stage 1} the new question \quad \textbf{stage 2} $80$ answers};
  \node[font=\scriptsize, anchor=west] at (5.5,0.55) {\textbf{stage 3} yes / no counts \quad \textbf{stage 4} the sentence};
\end{tikzpicture}\par}
\\
& \notesprompt{We work these together.}
\begin{probgrid}{|Y|Y|}
\pcell{16}{Name the variable of interest and its type --- \emph{before} looking at the $80$ answers. How do you know?}{1.2cm}{Would they ride --- yes or no. Categorical.} & \pcell{17}{What percent of the students asked said yes? If they are typical of the school, about how many of the $900$ is that?}{1.2cm}{$24/80 = 30\%$; $0.30 \times 900 = 270$.} \\ \hline
\end{probgrid}
\\
& \begin{probgrid}*{|Y|Y|}
\pcell{18}{Write the conclusion the council can hand to the district. Name the group asked, the number, and what was measured.}{1.6cm}{``Of the 80 Riverbend students asked, 30\% would ride if there were racks.''} & \pcell{19}{Suppose the $80$ had all been standing at the bike rack. Which stage would the problem have entered at --- and what could the report no longer claim?}{1.6cm}{Stage 2 --- only riders stand at the rack. It cannot describe all students.} \\ \hline
\end{probgrid}
\\ \hline
\end{guidednotes}

% ── THE NOTES END AT GUIDED PRACTICE ─────────────────────────────────────────
% There is deliberately no "you do" here: ../homework/ is the individual practice,
% started in class in the last ten minutes while the teacher circulates.

\end{document}
